\section{Quantum Foundations}

The base of this theory is \textbf{local}, \textbf{discrete}, and \textbf{deterministic}. The mesh is a fixed weave of docks, each dock carries $24$ sites, each site holds one tone in $\{-1,0,+1\}$, and the knit advances the whole configuration one beat at a time by a fixed reversible table. Nothing in that picture reaches across the mesh. A dock sees only its own sites and the sites that stream into it from its neighbours. Quantum mechanics, by contrast, is \textbf{nonlocal}. Two entangled measurements, far apart, show correlations that no local model can reproduce. That is a sharp tension, and it would be dishonest to bury it.

\begin{principle}[The tension, stated plainly]
A local discrete base cannot violate the Bell bound. Physics on the holographic boundary can. The base stays local, and nonlocality is an emergent boundary phenomenon.
\end{principle}

This section states the tension honestly, names exactly what is derived and what stays open, and shows where the substrate already supplies the missing ingredient. The resolution is not a patch. It is the same holographic geometry that the gravity and cosmology sections already need, doing one more job.

\subsection{The Bell tension}

A \emph{local hidden-variable theory} is one in which every outcome is fixed in advance by variables carried locally, with no faster-than-light coordination. The bare base is exactly such a theory. The knit is local and deterministic, so the tone configuration at any beat is a local hidden variable for everything that follows. Any theory of that form obeys the Bell bound, expressed through the CHSH quantity, which can never exceed $2$.

\begin{theorem}[Bell bound on the bare base]
The knit, being a local deterministic law over the mesh, is a local hidden-variable theory. Its CHSH correlator is therefore bounded by $2$. This bound is saturated but never exceeded by any deterministic strategy, a fact confirmed by brute-force search over all such strategies.
\end{theorem}

Quantum mechanics goes further. The optimal quantum measurement angles push the CHSH correlator up to the Tsirelson bound, $2\sqrt{2} \approx 2.83$. The two numbers are collected in Table~\ref{tab:bell-bounds}.

\begin{table}[ht]
\centering
\begin{tabular}{@{}lll@{}}
\toprule
theory & CHSH limit & status \\
\midrule
local hidden variable (the bare base) & $\le 2$ & verified by brute force over deterministic strategies \\
quantum mechanics & $2\sqrt{2} \approx 2.83$ & the Tsirelson bound, verified at the optimal angles \\
\bottomrule
\end{tabular}
\caption{The Bell gap. A local discrete base is bounded by $2$. Quantum correlations reach $2\sqrt{2}$.}
\label{tab:bell-bounds}
\end{table}

The gap is real. Two versus $2\sqrt{2}$ is not a rounding error, and it is not something a cleverer local rule can close.

\begin{principle}[No local escape]
A local discrete base, taken literally as the whole story, cannot reproduce quantum nonlocality. This is an honest tension, not a free pass, and we do not wave it away.
\end{principle}

We resolve it through the geometry the substrate already carries, not by smuggling nonlocality into the knit.

\subsection{The resolution: emergent holographic nonlocality}

The key move is to notice where physics actually lives. The bulk knit is local, so the \textbf{bulk} obeys $\mathrm{CHSH} \le 2$. But the observable physics of this theory does not live in the $4$-dimensional bulk. It lives on the holographic \textbf{boundary}, the cusp of the mesh, where the renormalization tree terminates.

The boundary inherits a very different notion of distance. Two points that sit far apart along the boundary can be close \emph{through the bulk}, linked by a short bulk-tree propagator that dives into the $4$-dimensional interior and surfaces elsewhere. Concretely:

\begin{itemize}
\item Two boundary points are connected \textbf{through the bulk}, by the bulk-tree propagator.
\item A single path through the bulk links boundary points that are widely separated along the boundary itself.
\item So the \textbf{boundary theory is nonlocal}. Its correlations are bulk-mediated, not boundary-local.
\end{itemize}

A bulk-mediated boundary theory is \textbf{not} a local hidden-variable model. The hidden variables of the boundary live in the bulk, and the bulk reaches across boundary separations. Such a theory is free to violate the Bell bound, up to the quantum value. The two layers are set side by side in Table~\ref{tab:two-layers}.

\begin{table}[ht]
\centering
\begin{tabular}{@{}lll@{}}
\toprule
layer & locality & Bell \\
\midrule
the bulk (the base) & local & $\le 2$ \\
the boundary (physics) & nonlocal, bulk-mediated & can reach $2\sqrt{2}$ \\
\bottomrule
\end{tabular}
\caption{Locality splits by layer. The base stays local and bounded. The emergent boundary is nonlocal and can reach the Tsirelson value.}
\label{tab:two-layers}
\end{table}

\begin{result}[Nonlocality is emergent]
Bell violation is an emergent boundary phenomenon, not a base one. The base stays local and discrete. The boundary, being bulk-mediated, can violate the Bell bound up to $2\sqrt{2}$, with no nonlocality ever inserted into the knit.
\end{result}

This is consistent with, and in fact requires, the holographic structure the substrate already has for independent reasons. The same bulk that mediates the gravity correlator (Section~\ref{sec:gravity}) and that branches into the renormalization tree (Section~\ref{sec:cosmology}) is what carries the nonlocal correlations here. One structure does all three jobs.

\begin{principle}[The bulk as the deep medium]
The bulk is the deep medium through which distant ripples on the surface stay in touch. Gravity, cosmic structure, and quantum nonlocality are three readings of the same bulk connectivity.
\end{principle}

\subsection{Where the complex \texorpdfstring{$i$}{i} comes from}

A natural worry remains even once nonlocality is handled. The base is \textbf{real and discrete}. Every tone is $-1$, $0$, or $+1$, and every beat is an integer step. Quantum amplitudes, however, are \textbf{complex}. Where does the imaginary unit $i$ come from, if the substrate never contains one?

There is a clean answer through the \emph{Doi-Peliti} map, the standard correspondence between a stochastic particle process and a second-quantized operator theory. The map applies because the knit's moves are exactly of the right algebraic form.

\begin{itemize}
\item The knit's three elementary moves are \textbf{create}, \textbf{annihilate}, and \textbf{hop}.
\item All three are the \textbf{same operator}, the transfer of one unit of $S^z$ across an edge between adjacent docks.
\item So the knit is exactly a \textbf{spin-$1$ XX model}, with the create move setting the source term.
\end{itemize}

The reversibility of the knit is what turns this stochastic model into a genuine quantum one. The knit collides then streams by a table that is reversible and charge-conserving, which is the lattice statement of detailed balance.

\begin{itemize}
\item The knit's reversibility (detailed balance) is precisely what makes the Doi-Peliti generator \textbf{Hermitian}.
\item A Hermitian generator is a real, bounded, positive-energy quantum Hamiltonian.
\item In the spin-coherent-state path integral of that model, the \textbf{Berry phase} term is the source of the quantum $i$.
\end{itemize}

\begin{result}[The complex unit is the Berry phase]
The complex amplitude is not added by hand. It is the Berry phase of the knit's own spin-exchange. The imaginary unit is where the discrete walk's phase lives, made visible by passing to the spin-coherent path integral of the spin-$1$ XX model the knit already is.
\end{result}

\subsection{The Born rule}

The $|\psi|^2$ measure is the natural weighting of the discrete walk's oscillation amplitudes. When the spin-exchange walk is read as a quantum amplitude, the amplitude-squared weighting is the measure of the underlying real walk, so the statistical form of the Born rule emerges rather than being imposed.

\begin{itemize}
\item The amplitude-squared weighting emerges as the \textbf{measure of the real walk}.
\item The full account of measurement and collapse is the \textbf{deeper open piece}.
\item This gap is shared with every discrete and emergent-quantum approach. No program has closed it.
\end{itemize}

We claim the measure, not the collapse. That is the honest line, and we draw it sharply rather than dressing a partial result as a complete one.

\subsection{Reflection positivity and entanglement}

For the emergent boundary theory to be a \textbf{genuine quantum theory}, and not merely a formal one, it must rest on reflection positivity.

\begin{definition}[Reflection positivity]
\emph{Reflection positivity} is the Euclidean condition guaranteeing a Hilbert space with positive norms and a unitary continuation to real time. A Euclidean theory that satisfies it defines a sensible quantum theory. One that fails it does not.
\end{definition}

The knit supplies exactly this condition through its reversibility.

\begin{itemize}
\item The knit's reversibility (detailed balance, the Hermitian Doi-Peliti generator) is precisely the property that yields reflection positivity.
\item A non-reversible knit would give a non-Hermitian generator, with negative norms and no consistent quantum theory.
\end{itemize}

So reversibility is not a cosmetic choice in the base. It is what makes the emergent theory quantum-mechanically sound.

\begin{lemma}[Reversibility gives the Hilbert space]
The reversibility of the knit makes the Doi-Peliti generator Hermitian, which supplies reflection positivity, which guarantees a positive-norm Hilbert space and unitary real-time evolution on the boundary.
\end{lemma}

Entanglement is likewise structural, not an added ingredient. The radial bulk tree is a \emph{MERA} network, the multiscale entanglement renormalization ansatz (Section~\ref{sec:cosmology}, with the holographic-code view in Section~\ref{sec:holography}).

\begin{itemize}
\item A MERA network builds boundary states with a definite entanglement structure, area-law entanglement entropy, by construction.
\item So boundary entanglement is the \textbf{default} of the holographic geometry, not a separate postulate.
\item This is the same tree that mediates the nonlocal Bell correlations. Entanglement and nonlocality are two faces of the bulk connectivity.
\end{itemize}

\subsection{Why a local base needs an emergent nonlocal boundary}

The pieces fit into a single ledger. Each requirement of a quantum theory is met by some feature the substrate already has, and the source of each is named in Table~\ref{tab:quantum-ledger}.

\begin{table}[ht]
\centering
\begin{tabular}{@{}ll@{}}
\toprule
requirement & supplied by \\
\midrule
a positive-norm Hilbert space & reflection positivity, from knit reversibility \\
complex amplitudes & the Berry phase of the spin-exchange \\
the $|\psi|^2$ measure & the discrete walk's oscillation amplitudes (collapse still open) \\
Bell violation up to $2\sqrt{2}$ & bulk-mediated boundary nonlocality \\
area-law entanglement & the MERA bulk tree \\
\bottomrule
\end{tabular}
\caption{The quantum ledger. Each ingredient of a quantum theory is supplied by a feature the substrate already carries.}
\label{tab:quantum-ledger}
\end{table}

A purely local discrete knit cannot be quantum on its own. It is bounded by $\mathrm{CHSH} \le 2$. What rescues it is \textbf{not} breaking locality in the base. It is the holographic geometry, which makes the \textbf{boundary} nonlocal while the \textbf{base} stays local.

\begin{principle}[The bridge]
The discrete local base needs an emergent nonlocal boundary to be quantum. The substrate supplies exactly that, through the same bulk that does gravity and cosmology. The base stays local, the boundary goes quantum, and the bulk is the bridge.
\end{principle}

\subsection{Connection to the monism}

The base of this theory is \textbf{experiential}. The tones are pain, peace, and pleasure, the values $-1$, $0$, and $+1$, and the substance carried by every site is a vibe. The substrate is mind-like at bottom, so the quantum amplitudes are the \textbf{texture of the underlying noticing}. The universe as a vibe mesh is where the quantum layer ultimately sits.

This is an honest \emph{pointer}, not a derivation. We do not claim to derive the amplitudes from the experiential character of the base, nor the experiential character from the amplitudes. We note the framing the monism makes available, and we stop there.

\subsection{Status, stated honestly}

The full accounting of what is settled and what is open is given in Table~\ref{tab:quantum-status}. The tension is real, the resolution is structural and rests on holography the substrate independently needs, and the deepest piece, the measurement problem, stays open and is not hidden.

\begin{table}[ht]
\centering
\begin{tabular}{@{}ll@{}}
\toprule
claim & status \\
\midrule
base obeys $\mathrm{CHSH} \le 2$ & derived, verified \\
quantum reaches $2\sqrt{2}$ & known (Tsirelson), verified \\
boundary can violate Bell via bulk mediation & derived (structural) \\
complex $i$ from the Berry phase & derived (Doi-Peliti) \\
reflection positivity from reversibility & derived \\
the $|\psi|^2$ measure & partial, measure emerges \\
collapse and measurement & open, shared with all discrete approaches \\
\bottomrule
\end{tabular}
\caption{Honest status of the quantum foundations. The structural claims are derived. The measurement problem stays open.}
\label{tab:quantum-status}
\end{table}

\begin{result}[The honest line]
The base stays local. The boundary goes quantum. The bulk is the bridge. The measure of the Born rule emerges, the collapse does not, and we say so rather than claiming a closure no discrete program has reached.
\end{result}

\subsection{See also}

Section~\ref{sec:holography} owns the Ryu-Takayanagi law, the area law, and the holographic codes the boundary nonlocality rides on. Section~\ref{sec:cosmology} carries the MERA bulk tree. Section~\ref{sec:gravity} carries the bulk-mediated correlator that shares this bulk.
